% DATED POST-REGISTRATION ADDENDUM: 2026-11-02.
% Theory-owned; do not edit any frozen v5 forecast or its comparison rule.
\subsection{Test 6: first prospective test on image-disjoint ImageNet images}
\label{app:s3-test6-xpn-20261102}

\paragraph{Registration and scoring.}
The v5 code was frozen before the new ImageNet images' outcomes
were available (code commit \texttt{81003cc}), and the numerical
v5 forecast was sealed and pushed on 2026-11-01 at 19:53:58 UTC
(XP commit \texttt{55ba86b}), before any run with index
$4$--$63$ was scored. The model was fitted only on runs
$0$--$3$ of the 100 \emph{image-disjoint} test images.
The registered primary comparison is race-model predictive-interval
coverage for E0 and E1 (class-Beta decision model), across
$K\in\{4,8,16,32,64\}$ from the \emph{complete}
64-run panel. Neither its intervals nor its model-selection
rule is revised in light of the results.

On the ten predeclared selection points, race covers
\textbf{10/10}, followed by the stronger-regularization
race variant \textbf{7/10}, the transferred tail variant
\textbf{5/10}, and exchangeable ranks S2 \textbf{0/10}.
Their mean absolute standardized errors are
$1.236$, $1.329$, $2.258$, and $7.441$, respectively.
Race is the winner under the frozen rule.
The five $K$-points share the \emph{same} images, candidate panels,
and nested subset evaluations; the ten coverage indicators are
correlated, not ten independent Bernoulli trials.
Combining the prospectively scored tests 2--6 gives
race \textbf{67/69} versus S2 \textbf{17/69}
selection-coverage points; this pooled descriptive tally
is likewise correlated and does not establish exact calibration
of a standalone generative mechanism.

\paragraph{Observed reconstruction and oracle curves.}
The complete-panel successes are
\[
\begin{array}{c|ccccc}
K&4&8&16&32&64\\ \hline
E0&.359&.415&.463&.501&.520\\
E1&.456&.498&.542&.585&.610\\
O_{\rm raw}&.470&.564&.634&.690&.730\\
O_{\rm orbit}&.569&.624&.669&.711&.750
\end{array}
\]
For race at $K=64$, the frozen E0 forecast is
$.462\ [ .378,.544]$ versus $.520$ observed,
and E1 is $.626\ [ .534,.711]$ versus $.610$.
By contrast, the correct-orbit oracle forecast
$.880\ [ .821,.931]$ misses the observed $.750$,
and the registered orbit-oracle interval misses at
\emph{all five} $K$-values (\textbf{0/5}).
The raw-GOOD oracle covers \textbf{3/5} values and
is \emph{underpredicted} at $K=4,8$:
$.417$ versus $.470$ and $.510$ versus $.564$.
The orbit-generation miss is therefore not a simple
uniform downward shift in all oracle success rates.
The correct-orbit oracle was optimistic in five
of the six prospectively reported tests.

\paragraph{Which image-level dependency was missed?}
Let $G_K$ mean that a uniform $K$-subset includes at
least one raw-GOOD run, and let $O_K$ mean that it includes
a correct-orbit but non-GOOD run. Then, exactly,
\begin{equation}\label{eq:s3-test6-orbit-only}
 O_{\rm orbit}(K)-O_{\rm raw}(K)
 =\Pr(O_K\cap G_K^c).
\end{equation}
At $K=64$, the registered v5 difference is
$.164$, but its observed value is \emph{only}
$.750-.730=.020$. Thus the model expected about
$16.4$ of 100 images to be rescued \emph{only} by
an orientation-orbit candidate, compared with
two actually observed. Independent inspection
of the authorized 64-run class strings gives
\textbf{67} images with both a GOOD and a non-GOOD
correct-orbit run, \textbf{6} with GOOD only,
\textbf{2} with orbit-only, and \textbf{25} with
neither. Across the 100 images, the respective
64-run GOOD and non-GOOD correct-orbit class counts
have descriptive Pearson correlation approximately
$+0.369$. In particular, $67/69$ images containing
an orbit-error run also contain a raw-GOOD run.
The low increment in
\eqref{eq:s3-test6-orbit-only} therefore reflects
\emph{co-occurrence across images}, not the absence
of orbit-error candidates.

The v5 hierarchical five-class Dirichlet model
and its observed-prefix conditioning did not
represent this image-level overlap adequately
when extrapolated to 64 runs. A single Dirichlet
distribution for a new image's class probabilities
also imposes negative covariance between its
distinct class-probability coordinates, whereas
the observed long-run class counts are positively
associated across images. This is evidence of
a shared easy/hard propensity or other dependence
outside that simple prior, although the observed
count correlation alone does not identify a unique
alternative joint mixing law or causal mechanism.
It is a \emph{model-misspecification diagnosis}:
the primary selector's accurate coverage can coexist
with offsetting errors in the oracle and
oracle-to-selector selection-loss models.

\paragraph{Relation to finite-moment non-identifiability.}
If a fixed image has per-run class probabilities
$\theta_G$ (raw GOOD) and
$\theta_O$ (other correct-orbit), then the
idealized population increment is the joint
mixing functional
\begin{equation}\label{eq:s3-test6-joint-functional}
 D_K=\mathbb E\!\left[
 (1-\theta_G)^K-
 (1-\theta_G-\theta_O)^K
 \right].
\end{equation}
Even the \emph{complete distribution} of
four-run five-class counts determines only
polynomial moments of total degree at most four
of the joint law of $(\theta_G,\theta_O,\ldots)$,
not the degree-$64$ functional
\eqref{eq:s3-test6-joint-functional}.
The C-T8 truncated-moment obstruction is
therefore a reason why correct long-run
coverage cannot be guaranteed from short-run
counts, and applies more broadly to the
\emph{joint} orbit-only event.
However, it neither predicts the sign nor
explains the magnitude of v5's $14.4$-point
orbit-only excess. The concrete, independently
released 64-run overlap table demonstrates
that the \emph{particular registered Dirichlet
extrapolation} fails on the true new-image
joint dependence. Non-identifiability is a
limitation of available calibration information;
misspecification describes this model's
actual out-of-sample prediction error.
Neither statement implies mathematical
unsolvability for an individual image:
$25/100$ images simply have no correct-orbit
candidate in the observed 64-run panel.

\paragraph{Separate stopping pre-score (not v5 fitting).}
The independently frozen stopping prediction
called N-H1a a pass at $N_{\max}=16$
(central gain $+0.2$ points, stated range
$[-1.8,+2.2]$), but the observed
equal-mean-runs gain was $-1.516$ points,
so the \emph{binary prospective call failed}
despite the observation lying inside its
informal range.
N-H1b at $N_{\max}=64$ was correctly called
to pass its negative-gain criterion
(predicted $-2.7$, observed $-3.102$ points).
Image-bootstrap intervals are
$[-3.60,+0.47]$ and $[-5.96,-0.10]$
points, respectively, on the same 100 images.
These stopping observations are distinct from
the v5 E0/E1 predictive-interval scoring,
and are not used to refit any forecast.
